% !TEX program = xelatex
% !TEX encoding = UTF-8
\documentclass{beamer}
\usepackage{amsfonts,amsmath}
\usetheme{sintef}
\usepackage{xeCJK}
\setCJKmainfont{Microsoft YaHei}[BoldFont=Microsoft YaHei Bold]
\setCJKsansfont{Microsoft YaHei}[BoldFont=Microsoft YaHei Bold]
\setCJKmonofont{Microsoft YaHei}[Scale=0.92]

\usefonttheme[onlymath]{serif}

\usepackage{booktabs}
\usepackage{multirow}
\usepackage{tikz}
\usetikzlibrary{arrows.meta, positioning}

\definecolor{TechDark}{HTML}{0f172a}
\definecolor{TechAccent}{HTML}{0891b2}

% 相对本 .tex 所在目录的显式路径（\IfFileExists 不搜索 \graphicspath）
\newcommand{\imgdata}[1]{../../mydata_png/#1}
\newcommand{\imgresult}[1]{../../results/#1}
\newcommand{\imgslide}[1]{../img/#1}

\newcommand{\maybeincludegraphics}[2][width=\linewidth]{%
  \IfFileExists{#2}{%
    \includegraphics[#1]{#2}%
  }{%
    \fbox{\parbox[c][2.2cm][c]{0.9\linewidth}{\centering\scriptsize 待插入图像\\ \path{#2}}}%
  }%
}

\titlebackground*{assets/background}

\title{RetinexFormer 复现与改进}
\subtitle{计算摄像学大作业汇报}
\course{北京大学 · 信科学院}
\author{潘嘉铭、胡林敏、石杨子然}
\date{\today}

\begin{document}
\maketitle

% ------------------------------------------------------------------ 2. 问题与任务
\begin{frame}{问题与任务}
  \textbf{背景：}低光照图像的传统增强算法中存在亮度不足、噪声放大，细节丢失等问题，论文《Retinexformer: One-stage Retinex-based
  Transformer for Low-light Image Enhancement》结合Retinex物理理论和Transformer模型，兼顾图片的提亮效果和保真度。

  \vspace{0.6em}
  \textbf{本作业工作：}
  \begin{enumerate}
    \item \textbf{复现} RetinexFormer，在 5 个公开基准数据集上验证
    \item \textbf{自采验证}：两类场景（极限暗光纹理 / 高反差复杂光源），DNG/JPG 实拍数据
    \item \textbf{改进探索}：传统增强 baseline；LoRA 参数高效微调（mydata 子集验证）
  \end{enumerate}
\end{frame}

% ------------------------------------------------------------------ 3. Retinex 理论
\begin{frame}{Retinex 理论}
  \begin{columns}[T]
    \column{0.52\textwidth}
    \textbf{核心分解：}
    \[
      I(x,y) = R(x,y) \odot L(x,y)
    \]
    \begin{itemize}
      \item $I$：观测图像
      \item $L$：光照（Illumination），缓慢变化
      \item $R$：反射率（Reflectance），承载纹理与颜色
    \end{itemize}
    \vspace{0.5em}
    \small 思路：先估计并校正光照，再恢复反射细节——RetinexFormer 将此过程神经网络化。

    \column{0.46\textwidth}
    \centering
    \begin{tikzpicture}[
      box/.style={draw, rounded corners, minimum width=2.6cm, minimum height=0.75cm, align=center, font=\small},
      arr/.style={-{Stealth[length=2mm]}, thick},
      lbl/.style={font=\scriptsize, fill=white, inner sep=1pt}
    ]
      \node[box, fill=blue!8] (I) at (0,2.8) {低光输入 $I$};
      \node[box, fill=yellow!15] (L) at (-1.8,1.4) {光照 $L$};
      \node[box, fill=green!10] (R) at (1.8,1.4) {反射 $R$};
      \node[box, fill=orange!12] (O) at (0,0) {增强输出};

      \draw[arr] (I) -- (L);
      \draw[arr] (I) -- (R);
      \draw[arr] (L) -- (O) node[midway, left=2pt, lbl] {提亮、去噪};
      \draw[arr] (R) -- (O);
    \end{tikzpicture}
  \end{columns}
\end{frame}

% ------------------------------------------------------------------ 3.5 RetinexFormer 理论重构
\begin{frame}[t]{理论重构：为真实物理世界的「不完美」建模}
  \vspace{-0.55em}
  \footnotesize
  \begin{columns}[T,onlytextwidth]
    \column{0.54\textwidth}

    \begin{block}{经典 Retinex 陷阱}
      \centering
      $I = R \odot L$
      \vspace{0.15em}

      \raggedright\scriptsize
      假设世界是完美无噪的，导致提亮后噪点爆炸。
    \end{block}

    \vspace{0.12em}
    \begin{block}{RetinexFormer 真实物理模型}
      \centering
      $I = (R + \hat{R}) \odot (L + \hat{L})$

      \vspace{0.1em}
      {\scriptsize $\Downarrow$ 引入提亮图 $\overline{L}$ 并化简}
      \[
        I_{lu} = I \odot \overline{L} = R + C
      \]

      \raggedright\scriptsize
      $C$：\textbf{整体退化项}（高 ISO 噪点 + 光照色偏）
    \end{block}

    \vspace{0.12em}
    \scriptsize
    引入物理扰动项，锁定暗光场景下放大后的高 ISO 噪点与光照色偏。

    \column{0.43\textwidth}
    \centering
    \begin{tikzpicture}[
      box/.style={draw, rounded corners, minimum width=2.9cm, minimum height=0.58cm, align=center, font=\scriptsize},
      arr/.style={-{Stealth[length=1.5mm]}, thick, TechAccent},
      lbl/.style={font=\tiny, fill=white, inner sep=1.5pt, text=TechDark}
    ]
      \node[box, fill=blue!8] (I) at (0,2.1) {低光输入 $I$};
      \node[box, fill=yellow!15] (Ilu) at (0,0.95) {$I_{lu} = I \odot \overline{L}$};
      \node[box, fill=green!10] (R) at (0,-0.15) {清晰反射 $R$};
      \draw[arr] (I) -- (Ilu);
      \draw[arr] (Ilu) -- (R);
      \node[lbl, anchor=west] at (0.12, 1.55) {Estimator};
      \node[lbl, anchor=west] at (0.12, 0.4) {Denoiser};
    \end{tikzpicture}

    \vspace{0.15em}
    \begin{block}{导向下一步架构}
      \scriptsize
      \begin{itemize}
        \setlength{\itemsep}{0pt}
        \setlength{\parsep}{0pt}
        \item \textbf{Estimator}：预测 $\overline{L}$，完成物理提亮
        \item \textbf{Denoiser}：拟合残差，去除 $C$ 并恢复 $R$
      \end{itemize}
    \end{block}
  \end{columns}
\end{frame}

% ------------------------------------------------------------------ 4. 模型结构
\begin{frame}{模型结构}
  \centering
  \begin{tikzpicture}[
    box/.style={draw, rounded corners, minimum height=0.85cm, align=center, font=\small, inner sep=4pt},
    arr/.style={-{Stealth[length=2mm]}, thick}
  ]
    \node[box, fill=blue!10, minimum width=2.2cm] (in) {低光输入};
    \node[box, fill=yellow!15, minimum width=2.8cm, right=0.55cm of in] (est) {Illumination\\Estimator};
    \node[box, fill=green!10, minimum width=2.6cm, right=0.55cm of est] (den) {Denoiser\\U-Net + IGAB(Transformer)};
    \node[box, fill=orange!12, minimum width=2.2cm, right=0.55cm of den] (out) {高清输出};

    \node[box, fill=yellow!8, minimum width=2.5cm, below=0.65cm of est] (fea) {\scriptsize $\overline{L}$ \;/\; $F_{lu}$};

    \draw[arr] (in) -- (est);
    \draw[arr] (est) -- (den);
    \draw[arr] (den) -- (out);
    \draw[arr] (est) -- (fea);
    \draw[arr] (fea) -- (den);

    \node[align=left, font=\scriptsize, below=0.35cm of fea] (igmsa) {
      \textbf{IG-MSA：}$F_{lu}$ 与 Value 哈达玛积 $\Rightarrow$ 暗区自适应去噪
    };
  \end{tikzpicture}

  \vspace{0.4em}
  \small
  Estimator：4 通道输入（RGB + 亮度均值）$\rightarrow$ $\overline{L}$（提亮图）与 $F_{lu}$（光照特征）；\\
  Denoiser：Encoder--Decoder + 跳跃连接，修复 $\overline{L}$ 提亮带来的噪声与伪影。\\
  \scriptsize IGAB 中 FeedForward 采用 $1{\times}1$ + $3{\times}3$ 深度可分离卷积，增强局部连续性与边缘保持。
\end{frame}

% ------------------------------------------------------------------ 5. 公开数据集复现与 Self-ensemble
\begin{frame}{公开数据集复现与 Self-ensemble}
  \begin{columns}[T]
    \column{0.58\textwidth}
    \scriptsize
    \renewcommand{\arraystretch}{1.15}
    \begin{tabular}{lcccc}
      \toprule
      \multirow{2}{*}{数据集} & \multicolumn{2}{c}{无 SE} & \multicolumn{2}{c}{有 SE} \\
      \cmidrule(lr){2-3}\cmidrule(lr){4-5}
      & PSNR & SSIM & PSNR & SSIM \\
      \midrule
      LOL\_V1      & 25.15 & 0.845 & 25.27 & 0.853 \\
      LOL\_V2\_real & 22.79 & 0.839 & 22.82 & 0.850 \\
      LOL\_V2\_syn  & 25.67 & 0.930 & 25.83 & 0.934 \\
      FiveK        & 24.94 & 0.907 & 24.97 & 0.909 \\
      SID          & 24.44 & 0.680 & 24.36 & 0.681 \\
      \bottomrule
    \end{tabular}
    \tiny SE = Self-ensemble；5 个公开集复现成功。\\
    SMID / SDSD / NTIRE 2024 因数据体量、文件损坏或链接失效未纳入公开集复现。

    \column{0.40\textwidth}
    \small
    \textbf{Self-ensemble 策略}
    \begin{itemize}
      \item 测试时对输入做 8 种几何变换（翻转 / 旋转组合）
      \item 分别推理后逆变换，对 8 路输出取平均
      \item 以约 $8\times$ 推理时间换取更稳定结果
    \end{itemize}
    \vspace{0.3em}
    \scriptsize
    \textbf{提升小结：}LOL 系列 SSIM +0.005$\sim$0.01；\\
    高分辨率集收益较小，但无明显负效应
  \end{columns}
\end{frame}

% ------------------------------------------------------------------ 6. 自采数据：流程与结果
\begin{frame}{自采数据：流程与结果}
  \vspace{-0.65em}
  \begin{columns}[T]
    \column{0.38\textwidth}
    \scriptsize
    \textbf{数据采集}
    \begin{tabular}{lcc}
      \toprule
      场景 & 有 GT & 无 GT \\
      \midrule
      极限暗光纹理   & 9 对 & 4 张 \\
      高反差复杂光源 & 5 对 & --- \\
      \bottomrule
    \end{tabular}

    \vspace{0.08em}
    \textbf{流程：}\ {\tiny 实拍（DNG/JPG）$\rightarrow$ 裁剪/缩放 $\rightarrow$ PNG $\rightarrow$ 推理 $\rightarrow$ 评测}

    \column{0.58\textwidth}
    \centering
    \scriptsize
    \renewcommand{\arraystretch}{1.12}
    \begin{tabular}{lcc}
      \toprule
      场景 & LOL\_v1 & LOL\_v2\_real \\
      \midrule
      极限暗光 &
        15.01 (0.642) & 15.85 (0.670) \\
      复杂光源 &
        16.59 (0.662) & 14.91 (0.701) \\
      \bottomrule
    \end{tabular}\\[0.2em]
    {\scriptsize JPG，无 SE；PSNR (SSIM)，各场景 GT 对均值}
  \end{columns}

  \vspace{0.1em}
  \begin{columns}[T]
    \column{0.495\textwidth}
    \centering\textbf{\scriptsize 极限暗光 \#6}
    \column{0.495\textwidth}
    \centering\textbf{\scriptsize 复杂光源 \#4}
  \end{columns}
  \vspace{-0.05em}
  \begin{columns}[T]
    \column{0.165\textwidth}
    \centering
    \maybeincludegraphics[width=\linewidth,height=2.35cm,keepaspectratio]{\imgdata{jpg/with_targets/extreme_lowlight_texture/input/6.png}}\\[-0.05em]
    \scriptsize 输入
    \column{0.165\textwidth}
    \centering
    \maybeincludegraphics[width=\linewidth,height=2.35cm,keepaspectratio]{\imgresult{mydata_png/jpg/with_targets/extreme_lowlight_texture/input/6.png}}\\[-0.05em]
    \scriptsize 输出
    \column{0.165\textwidth}
    \centering
    \maybeincludegraphics[width=\linewidth,height=2.35cm,keepaspectratio]{\imgdata{jpg/with_targets/extreme_lowlight_texture/target/6.png}}\\[-0.05em]
    \scriptsize GT
    \column{0.165\textwidth}
    \centering
    \maybeincludegraphics[width=\linewidth,height=2.35cm,keepaspectratio]{\imgdata{jpg/with_targets/high_contrast_complex_light/input/4.png}}\\[-0.05em]
    \scriptsize 输入
    \column{0.165\textwidth}
    \centering
    \maybeincludegraphics[width=\linewidth,height=2.35cm,keepaspectratio]{\imgresult{mydata_png/jpg/with_targets/high_contrast_complex_light/input/4.png}}\\[-0.05em]
    \scriptsize 输出
    \column{0.165\textwidth}
    \centering
    \maybeincludegraphics[width=\linewidth,height=2.35cm,keepaspectratio]{\imgdata{jpg/with_targets/high_contrast_complex_light/target/4.png}}\\[-0.05em]
    \scriptsize GT
  \end{columns}
\end{frame}

% ------------------------------------------------------------------ 7. RetinexFormer 结构（论文原图）
\begin{frame}[plain]
  \centering
  \maybeincludegraphics[width=\textwidth,height=0.92\paperheight,keepaspectratio]{\imgslide{structure.png}}
\end{frame}

% ------------------------------------------------------------------ 8. 改进：Domain Gap 与 LoRA 微调
\begin{frame}{改进：Domain Gap 与 LoRA 微调}
  \vspace{-0.35em}
  \begin{columns}[T,onlytextwidth]
    \column{0.50\textwidth}
    \scriptsize
    \textbf{动机：Domain Gap}
    \begin{itemize}
      \setlength{\itemsep}{0.05em}
      \setlength{\parsep}{0pt}
      \item 同一 mydata 换 checkpoint：FiveK SSIM $\approx$ 0.21 vs SDSD\_outdoor $\approx$ 0.65
      \item LOL\_v1 在 mydata：PSNR $\approx$ 15.8 dB（公开集 25+ dB）
      \item 传统 baseline（Gamma/CLAHE）已实现，仍难消除彩噪
    \end{itemize}

    \vspace{0.08em}
    \textbf{LoRA 微调}
    \begin{itemize}
      \setlength{\itemsep}{0.05em}
      \setlength{\parsep}{0pt}
      \item 冻结 Estimator + Denoiser 主干，仅训练 \texttt{lora\_} 旁路（$r{=}4$）
      \item 注入：IG-MSA $Q/K/V/\mathrm{Proj}$ + IGAB FFN $1{\times}1$ Conv
      \item 损失：L1 + $0.2(1{-}\mathrm{SSIM})$；训练集为 mydata 部分有 GT 样本
      \item mydata JPG 评测：SSIM 和 PSNR 均有明显提升。
    \end{itemize}

    \column{0.47\textwidth}
    \scriptsize
    \textbf{训练曲线与验证指标}

    \begin{center}
      \maybeincludegraphics[width=0.72\linewidth]{\imgslide{loss_I-pix_with_mydata.png}}\\[-0.1em]
      {\tiny loss\_I-pix（mydata）}

      \vspace{0.12em}
      \maybeincludegraphics[width=0.72\linewidth]{\imgslide{loss_I-ssim_with_mydata.png}}\\[-0.1em]
      {\tiny loss\_I-ssim（mydata）}
    \end{center}

    \vspace{0.05em}
    \renewcommand{\arraystretch}{1.05}
    \begin{tabular}{lcc}
      \toprule
      设置 & PSNR & SSIM \\
      \midrule
      预训练 & 14.52 & 0.606 \\
      LoRA 微调 & 18.15  & 0.701 \\
      \bottomrule
    \end{tabular}
  \end{columns}
\end{frame}

% ------------------------------------------------------------------ 9. 微调效果对比
\begin{frame}{微调效果对比}
  \vspace{-0.45em}
  \scriptsize
  自采 JPG：预训练 \textbf{LOL\_v2\_real} vs \textbf{LoRA 微调}（mydata 子集，$r{=}4$）

  \vspace{0.12em}
  \begin{columns}[T]
    \column{0.24\textwidth}
    \centering\footnotesize\textbf{输入}
    \column{0.24\textwidth}
    \centering\footnotesize\textbf{LOL\_v2\_real}
    \column{0.24\textwidth}
    \centering\footnotesize\textbf{LoRA 微调}
    \column{0.24\textwidth}
    \centering\footnotesize\textbf{GT}
  \end{columns}

  \vspace{0.08em}
  \begin{columns}[T]
    \column{0.24\textwidth}
    \centering
    \maybeincludegraphics[width=\linewidth,height=2.55cm,keepaspectratio]{\imgdata{jpg/with_targets/extreme_lowlight_texture/input/6.png}}
    \column{0.24\textwidth}
    \centering
    \maybeincludegraphics[width=\linewidth,height=2.55cm,keepaspectratio]{\imgresult{mydata_png_lolv2_real/dng/with_targets/extreme_lowlight_texture/input/6.png}}
    \column{0.24\textwidth}
    \centering
    \maybeincludegraphics[width=\linewidth,height=2.55cm,keepaspectratio]{\imgresult{mydata/6.png}}
    \column{0.24\textwidth}
    \centering
    \maybeincludegraphics[width=\linewidth,height=2.55cm,keepaspectratio]{\imgdata{jpg/with_targets/extreme_lowlight_texture/target/6.png}}
  \end{columns}
  \vspace{-0.15em}
  {\tiny 极限暗光纹理 \#6}

  \vspace{0.08em}
  \begin{columns}[T]
    \column{0.24\textwidth}
    \centering
    \maybeincludegraphics[width=\linewidth,height=2.55cm,keepaspectratio]{\imgdata{dng/with_targets/extreme_lowlight_texture/input/5.png}}
    \column{0.24\textwidth}
    \centering
    \maybeincludegraphics[width=\linewidth,height=2.55cm,keepaspectratio]{\imgresult{mydata_png_lolv2_real/dng/with_targets/extreme_lowlight_texture/input/5.png}}
    \column{0.24\textwidth}
    \centering
    \maybeincludegraphics[width=\linewidth,height=2.55cm,keepaspectratio]{\imgresult{mydata/5.png}}
    \column{0.24\textwidth}
    \centering
    \maybeincludegraphics[width=\linewidth,height=2.55cm,keepaspectratio]{\imgdata{dng/with_targets/extreme_lowlight_texture/target/5.png}}
  \end{columns}
  \vspace{-0.15em}
  {\tiny 高反差复杂光源 \#5}

  \vspace{0.05em}
  \tiny
  mydata 子集 LoRA 微调后两类场景 SSIM 均有提升，复杂光源更贴近 GT。
\end{frame}

% ------------------------------------------------------------------ 10. 总结与分工
\begin{frame}{总结与分工}
  \small
  \textbf{工作总结}
  \begin{itemize}
    \item 成功复现 RetinexFormer，5 个公开集指标与论文趋势一致
    \item 搭建 DNG/JPG 自采 pipeline，揭示预训练权重的 domain gap
    \item 实现 Gamma / CLAHE 传统 baseline，明确深度模型优势与改进方向
    \item 实现 LoRA 微调并在 mydata 子集上验证
  \end{itemize}

  \vspace{0.5em}
  \textbf{分工}
  \begin{tabular}{llc}
    \toprule
    成员 & 主要负责 & 贡献比例 \\
    \midrule
    潘嘉铭 & 基础验证、LoRA微调模型、损失函数改进 & 40\% \\
    石杨子然 & 数据预处理脚本、实拍测试数据集、期末汇报PPT & 30\% \\
    胡林敏 & 传统 baseline、GitHub仓库、期末总报告 & 30\% \\
    \bottomrule
  \end{tabular}
\end{frame}

% ------------------------------------------------------------------ 11. 感谢倾听
\begin{chapter}[assets/background_negative]{maincolor}{感谢倾听}
  \centering
  \vspace{1em}
  {\Large 欢迎提问与讨论}
\end{chapter}

\end{document}
